%==========================================================================
\section{Reference design}
Java's reflection (\code{java.lang.reflect}) is the closest analogue and the
design we follow most closely. We chose it because LO's object model is
essentially Java's: statically typed classes, single inheritance, vtable
dispatch and checked casts. Like Java, we split reflection into runtime class
metadata (\code{Class}), field access (\code{Field}) and method invocation
(\code{Method}). C\#'s reflection is nearly the same design.

We diverge from Java in six ways. \textbf{Return typing:} Java's
\code{Field.get} and \code{Method.invoke} return a boxed
\code{java.lang.Object} that the caller must cast (\code{Field} also has typed
accessors such as \code{getInt} and \code{getBoolean}); LO-5 names the
expected type T at the call site, as in \textit{getfield(T, e, n)} and
\textit{invoke(T, e, n, \ldots)}. \textbf{Type matching:} where Java applies
boxing, unboxing and widening primitive conversions, LO-5 requires an exact
match for primitive fields (int, bool, String); for a class-typed field, a
read accepts a field declared T or a subclass of T, and a write checks the
value's runtime class against the field's declared type, like Java's
array-store check. \textbf{Metadata queries:} Java returns heap-allocated
arrays (e.g.\ \code{Class.getFields()}); LO-5 uses the indexed ReflectInfo
queries (\S4) that return ints or Strings, such as \textit{field\_name(int i)}. \textbf{Overloading:}
Java's method lookup accounts for parameter signatures; LO has no
overloading, so methods are looked up by their unique names.
\textbf{Construction:} like Java's \code{Constructor.newInstance},
\textit{new\_instance(T, name, args\ldots)} creates an object reflectively,
from a class named by a String that must be T or a subclass of T, checked at
runtime; \textit{allocate(T, name)} creates an instance without running a
constructor, its fields taking their defaults, the closest analogue to how
Java serialization creates objects. \textbf{Field access:} Java's
\code{Class.getField(name)} finds only public fields, and private ones need
\code{getDeclaredField} plus \code{setAccessible(true)}; LO has no access
modifiers, so every user-declared field, inherited ones included, is
reachable by name.
